%=====================================================================
% PART VI OPENER
%=====================================================================
\thispagestyle{plain}
\structuralfigures{P6.}

\begin{center}
\begin{tcolorbox}[enhanced, width=0.92\textwidth, colback=goldhl!5,
  colframe=goldhl, boxrule=0pt, leftrule=5pt, arc=3pt, top=8pt, bottom=8pt]
\large\itshape\color{goldhl!45!black}
The book closes where the analysis runs out. Chapter 21 is one more well-solved problem; Chapter 22 is the class of problems nobody can solve efficiently and nobody can prove cannot be --- the largest open question in the subject, stated precisely enough to recognise.
\end{tcolorbox}
\end{center}

\vspace{6pt}

\dsfigh{%
\begin{tikzpicture}[
  cbox/.style={rounded corners=3pt, draw=goldhl, line width=1pt, fill=goldhl!7,
               text=goldhl!45!black, font=\small\bfseries, align=center,
               minimum width=3.5cm, minimum height=1.0cm},
  hbox/.style={rounded corners=3pt, draw=goldhl, line width=1.8pt, fill=goldhl!16,
               text=goldhl!45!black, font=\small\bfseries, align=center,
               minimum width=3.5cm, minimum height=1.0cm},
  arr/.style={-{Stealth[length=2.4mm]}, goldhl!60, line width=1pt},
  note/.style={font=\scriptsize\itshape, text=goldhl!45!black, align=center,
               fill=white, inner sep=2pt}
]
  \node[cbox] (c0) at (0.00,0) {Ch.\ 21\\String Matching};
  \node[hbox] (c1) at (4.30,0) {Ch.\ 22\\P and NP};
  \node[cbox] (c2) at (8.60,0) {Ch.\ 23\\Your Project};
  \draw[arr] (c0) -- (c1);
  \draw[arr] (c1) -- (c2);
  \node[note] at (2.15,-1.0) {where it runs out};
  \node[note] at (6.45,-1.0) {all of it, by you};
  \node[font=\scriptsize\itshape, text=accent, align=center] at (4.30,1.15)
       {the open question, stated precisely};
\end{tikzpicture}}%
{Reading path through Part~VI. Chapter 22 is the only chapter whose central claim is that nobody knows; Chapter 23 is where you use the rest.}{part6map}

\newpage
